A coleta de informações de um usuário pode ser \emph{explícita} ou \emph{implícita}\cite{viviani10}. Na coleta \emph{explícita}, é requisitado ao usuário que preencha informações através de um formulário pré-definido, enquanto na coleta \emph{implícita}, a informação é inferida a partir do estudo do comportamento dos usuários.

Informações a respeito de um usuário colhidas por uma aplicação podem ser úteis para outras aplicações que tenham o mesmo usuário \cite{viviani10}. No entanto, os diferentes \emph{modelos de usuário} utilizados pelas aplicações tornam difícil a utilização de informações de um mesmo \emph{perfil de usuário} de uma aplicação por outra. Além disso, algumas das informações do usuário podem ser válidas apenas dentro de um \emph{contexto} específico, o que dificulta ainda mais que diferentes aplicações compartilharem o mesmo perfil.

Visando resolver estas questões, o conceito de \emph{personalização inter-sistemas} foi introduzido \cite{niederee04} como forma de compartilhar informações de usuários entre diferentes aplicações. Como o maior problema deste tipo de modelo é o alinhamento entre as diferentes abordagens de personalização, o estudo de \emph{sistemas de modelos de usuário} é o principal desafio da pesquisa atual na área de personalização.\cite{viviani10}

O IEEE Standard Computer Dictionary \cite{ieeedict} define interoperabilidade como ``a capacidade de dois ou mais sistemas ou componentes trocarem informações e utilizar a informação que foi trocada.'' No âmbito da personalização, interoperabilidade é a capacidade de acessar e interpretar informação derivada de múltiplas fontes heterogênias e integrar esta informação em um modelo de usuário de granularidade própria\cite{carmagnola09}. Os dois principais modelos de interoperabilidade são: (i) \emph{interoperabilidade sintática}, onde as diferenças entre os modelos de usuário são resolvidos em nível de aplicação, e (ii) \emph{interoperabilidade semântica}, onde as diferenças entre os modelos são resolvidos em nível de conhecimento.

As duas principais abordagens utilizadas em sistemas de modelos de usuários interoperáveis são \cite{viviani10}:

\begin{description}

  \item[Baseados em padrões]: são baseados numa visão \emph{top-down}, definindo um padrão que todas as aplicações envolvidas devem observar;
  \item[Baseados em mediação]: se baseiam numa visão \emph{bottom-up}, onde ocorre algum tipo de \emph{reconciliação} entre os diferentes modelos de representação.

\end{description}

\todo[inline]{Descrever conceitos de perfis relacionados a ontologias.}

%
% Perfis
% Computação Ubíqua
% Trilhas


